\chapter{Divisibility and factorization in domains}\label{chap:M07.10}

Throughout this chapter a domain is commutative and nontrivial.
% Topic grouping: Divisibility and cancellation
\KEntry{K07.10.divisibility}{Divisibility}
Write $a\mid b$ if $b=ac$ for some $c\in R$. Then $a\mid0$ for every
$a$, and $0\mid b$ exactly when $b=0$. Divisibility is reflexive and
transitive, since $a=a1$ and $b=ac$, $d=be$ imply $d=a(ce)$.
A unit divides every element. Conversely an element dividing $1$ is a
unit. The relation depends on the ring; for example $2\mid1$ in $\Q$
and not in $\Z$.

\KEntry{K07.10.associates}{Associates}
Elements $a,b$ are associates if $a=ub$ for a unit $u$.
This is an equivalence relation: use $1$, $u^{-1}$ and products of units
for reflexivity, symmetry and transitivity. Associates divide each other.
Conversely, if $a=bc$ and $b=ad$ and $a\ne0$, then
$a=acd$ gives $cd=1$, so $c$ is a unit. If $a=0$, mutual divisibility
forces $b=0$. Thus
\[
 a\text{ and }b\text{ are associates}
 \quad\Longleftrightarrow\quad a\mid b,\ b\mid a
 \quad\Longleftrightarrow\quad (a)=(b).
\]
In $\Z$ the ambiguity is sign; in $F[X]$ it is multiplication by a
nonzero constant.

\KEntry{K07.10.gcd_lcm}{Greatest common divisors and least common multiples}
A gcd $d$ of $a,b$ divides both and is divisible by every common
divisor. An lcm is divisible by both and divides every common multiple.
If two gcds exist each divides the other, so they are associates;
the same argument applies to lcms. These definitions are divisibility
conditions, not assertions about a numerical order on $R$.
We allow the normalized conventions $\gcd(0,0)=0$ and
$\lcm(a,0)=0$. For $a\ne0$, a gcd of $a,0$ is $a$ up to units.
Not every domain has gcds for every pair.

\KEntry{K07.10.prime_irreducible}{Prime and irreducible elements}
An irreducible is a nonzero nonunit $p$ such that $p=ab$ forces $a$ or
$b$ to be a unit. A prime element is a nonzero nonunit $p$ with
$p\mid ab\Rightarrow p\mid a$ or $p\mid b$.
The latter condition is equivalent to primality of $(p)$, by the
principal-ideal divisibility criterion. The nonzero requirement for an
element is retained even though the zero \emph{ideal} can be prime.
Aluffi includes zero among prime elements; his convention must be
restricted to nonzero elements before comparing the next theorem.

\KEntry{K07.10.prime_implies_irreducible}{Prime implies irreducible}
\begin{proposition}\label{prop:prime-implies-irreducible}
A prime element of a domain is irreducible.
\end{proposition}
\begin{proof}
If $p=ab$, primality gives, say, $a=pc$. Then $p=pcb$, and
cancellation of $p\ne0$ gives $cb=1$, making $b$ a unit.
The other alternative gives a unit $a$.
\end{proof}

\KEntry{K07.10.irreducible_nonprime}{An irreducible which is not prime}
In $R=\Z[\sqrt{-5}]$, complex conjugation gives the multiplicative
integer norm $N(a+b\sqrt{-5})=a^2+5b^2$. Its value is $1$ precisely
at $1,-1$, which are therefore the units: a unit must have norm a
positive integer dividing $1$, and norm $1$ gives inverse by conjugation.
If $3=uv$ with nonunits, the positive norms multiply to $9$ and each
exceeds $1$, so both would be $3$. But $a^2+5b^2=3$ has no integral
solution: $b\ne0$ makes the left side at least $5$, while $b=0$ would
require $a^2=3$. Thus $3$ is irreducible. Nevertheless
\[
 (2+\sqrt{-5})(2-\sqrt{-5})=9,
\]
and neither factor is divisible by $3$, since divisibility by the
integer $3$ requires both integer coordinates to be divisible by $3$.
So $3$ is not prime.

% Topic grouping: Euclidean division and principal ideals
\KEntry{K07.10.euclidean_definition}{Euclidean domains}
A Euclidean domain has a well-founded size relation on remainders with
the property that for $b\ne0$ there exist $q,r$ with
$a=qb+r$ and either $r=0$ or $r$ smaller than $b$.
A function to $\N$ with strict decrease is a familiar sufficient
description. The relation, quotient and remainder data are additional
structure; a class name alone is not an algorithm or a complexity bound.
The integers are Euclidean using absolute value and ordinary division.
A field is Euclidean with every nonzero division having remainder zero.

\KEntry{K07.10.euclidean_algorithm}{The extended Euclidean algorithm}
\begin{theorem}\label{thm:euclidean-algorithm}
Repeated Euclidean division terminates. Its last nonzero remainder is
a gcd of the two inputs, and is an $R$-linear combination of them.
\end{theorem}
\begin{proof}
Start with $r_{-1}=a,r_0=b$, and for a nonzero current divisor write
$r_{i-1}=q_{i+1}r_i+r_{i+1}$. Nonzero successive remainders strictly
decrease, so well-foundedness forbids an infinite run. If the last
nonzero remainder is $d=r_k$, then $r_{k-1}=q_{k+1}d$.
Backward substitution shows $d$ divides every preceding remainder,
including $a,b$. Conversely a common divisor of $a,b$ divides
$r_{i+1}=r_{i-1}-q_{i+1}r_i$ inductively, hence divides $d$.
The same recurrence, starting with coefficient pairs $(1,0),(0,1)$,
expresses every $r_i$ as $x_i a+y_i b$.
If $b=0$, take $d=a$; if both inputs vanish, take $d=0$.
\end{proof}
Quotients are not generally unique for a chosen Euclidean function.
For example $5=2\cdot2+1=3\cdot2-1$ in $\Z$ has two remainders of
absolute value smaller than $2$. A chosen nonnegative remainder convention
selects one; the abstract definition does not.

\KEntry{K07.10.pid_definition}{Principal ideal domains}
A PID is a domain in which every ideal has one generator.
A principal ideal ring need not be a domain: the ideals of $\Z/4\Z$
are $0,(2)$ and the whole ring, but $2\cdot2=0$ there.
The ideal list follows from the correspondence with ideals of $\Z$
containing $(4)$, whose nonnegative generators are the divisors of $4$.

\KEntry{K07.10.euclidean_pid}{Euclidean implies PID}
\begin{theorem}\label{thm:euclidean-implies-pid}
Every Euclidean domain is a PID.
\end{theorem}
\begin{proof}
The zero ideal is generated by zero. In a nonzero ideal $I$, choose a
nonzero element $d$ minimal for the well-founded size relation restricted
to $I\setminus\{0\}$. For $a\in I$, divide $a=qd+r$.
The remainder belongs to $I$ and if nonzero would be smaller than $d$.
Thus $r=0$ and $a\in(d)$. The reverse inclusion follows from $d\in I$.
\end{proof}

\KEntry{K07.10.bezout}{Bezout identities}
If $(a,b)=(d)$, then $d$ divides $a,b$, and any common divisor $c$
has $(a,b)\subseteq(c)$, hence divides $d$. Thus $d$ is a gcd.
Also $d\in(a,b)$ means $d=xa+yb$.
This proves gcd existence and Bezout identities in a PID, with
the all-zero case $d=0$ included.

The converse assertion that gcd $1$ always yields a Bezout identity is
false. In $\Z[X]$, a divisor of $2$ must be a constant polynomial,
by additivity of degrees in a domain; a common divisor of $2,X$ must
then divide the coefficient $1$ of $X$, so is a unit. Hence their gcd
is $1$. But $2f+Xg$ always has even constant coefficient and cannot
equal $1$. The degree argument used here is proved in
Proposition~\ref{prop:polynomial-degree}.

\KEntry{K07.10.pid_prime_maximal}{Primes and irreducibles in a PID}
\begin{proposition}\label{prop:pid-primes}
In a PID every nonzero prime ideal is maximal. For a nonzero element
$p$, being irreducible, being prime, and having $(p)$ maximal are equivalent.
\end{proposition}
\begin{proof}
If $(p)$ is a nonzero prime ideal and $(p)\subseteq(d)$, write $p=dc$.
Primality gives $d\in(p)$, making the ideals equal, or $c=pe$, making
$p=dpe$ and hence $de=1$. In the latter case $(d)=R$.
For irreducible $p$, any containing ideal $(d)$ again gives $p=dc$,
so $d$ or $c$ is a unit. The former gives $(d)=R$ and the latter
$(d)=(p)$; thus $(p)$ is maximal.
Maximal ideals are prime by Proposition~\ref{prop:maximal-quotient},
and prime elements are irreducible by
Proposition~\ref{prop:prime-implies-irreducible}.
\end{proof}
The nonzero restriction on the ideal is needed: $(0)$ in $\Z$ is prime
and not maximal.

\begin{lemma}\label{lem:pid-acc}
Ascending chains of ideals in a PID stabilize. Every nonempty collection
of ideals has a maximal member.
\end{lemma}
\begin{proof}
For $I_1\subseteq I_2\subseteq\cdots$, the union is an ideal $(a)$.
Its generator lies in some $I_n$, forcing all later ideals to equal
the union. If a nonempty collection had no maximal member, successive
choices of strictly larger members would produce a nonstabilizing
ascending chain, a contradiction.
\end{proof}

% Topic grouping: Unique factorization
\KEntry{K07.10.ufd_definition}{The factorization requirement}
A UFD is a domain in which every nonzero nonunit has a finite factorization
into irreducibles, unique up to reordering and replacing corresponding
factors by associates. Equivalently one may write a leading unit times
chosen representatives of irreducibles. Zero is not given a prime
factorization. Units correspond to empty lists of nonunit factors.
A field is a UFD since every nonzero element is a unit.

\begin{lemma}\label{lem:prime-factorization-unique}
In a domain, factorizations into prime elements are unique up to units
and order.
\end{lemma}
\begin{proof}
If $u p_1\cdots p_m=v q_1\cdots q_n$, primality of $p_1$ makes it
divide a factor on the right: apply the binary prime condition repeatedly,
noting a nonunit prime cannot divide the unit $v$.
Say $p_1\mid q_j$. Irreducibility of $q_j$ makes the quotient a unit,
so $q_j$ is associate to $p_1$. Move that factor first and cancel the
nonzero prime, absorbing its unit into $v$. Induct on $m$.
If $m=0$, a nonempty product on the other side cannot be a unit:
if a product is a unit each factor divides $1$ and is a unit.
Thus $n=0$, completing the induction and proving equality of lengths.
\end{proof}

\KEntry{K07.10.pid_ufd}{PID implies UFD}
\begin{theorem}\label{thm:pid-implies-ufd}
Every PID is a UFD. Consequently every Euclidean domain is a UFD.
\end{theorem}
\begin{proof}
Suppose a nonzero nonunit lacks an irreducible factorization.
Among the principal ideals generated by such elements choose a maximal
one, using Lemma~\ref{lem:pid-acc}, say $(a)$.
The element $a$ is not irreducible, so $a=bc$ with both factors nonunits.
Both are nonzero. The inclusions $(a)\subsetneq(b),(c)$ are strict:
equality with $(b)$ would make $a,b$ associates and cancellation would
make $c$ a unit. By maximality both $b,c$ have irreducible
factorizations, giving one for $a$, a contradiction.
Thus existence holds. Proposition~\ref{prop:pid-primes} makes these
irreducibles prime, and Lemma~\ref{lem:prime-factorization-unique} proves
uniqueness.
\end{proof}

\KEntry{K07.10.ufd_prime_irreducible}{Irreducibles are prime in a UFD}
Let $p$ be irreducible in a UFD and suppose $ab=pc$.
If $a$ or $b$ is zero then $p$ divides that factor. Otherwise all
nonzero factors in the equation have irreducible factorizations
(units contribute empty lists). Uniqueness says $p$ is associate to a
factor occurring in $a$ or $b$, so it divides that factor.
Thus $p$ is prime. The reverse implication was already proved, giving
prime if and only if irreducible in a UFD.

\KEntry{K07.10.ufd_gcd}{Prime-exponent arithmetic}
Choose one representative for each associate class of primes appearing
in nonzero $a,b$, and write
$a=u\prod p^{e_p}$, $b=v\prod p^{f_p}$, using exponent zero when absent.
The element $\prod p^{\min(e_p,f_p)}$ is a gcd and
$\prod p^{\max(e_p,f_p)}$ an lcm. Indeed divisibility of nonzero
elements is exactly coordinatewise comparison of prime exponents,
as follows by factorizing the quotient and using uniqueness.
These formulas verify both universal divisibility conditions.
Since $\min(e,f)+\max(e,f)=e+f$, the product of gcd and lcm is
associated to $ab$. If an input is zero, the conventions stated earlier
give the same association, both products being zero. Literal equality
requires compatible choices of units; association does not.

\KEntry{K07.10.nonufd_domain}{Failure of unique factorization}
In $\Z[\sqrt{-5}]$ the identity
\[
 6=2\cdot3=(1+\sqrt{-5})(1-\sqrt{-5})
\]
gives two irreducible factorizations. For $2$, a nonunit factorization
would require norm $2$ factors, impossible by $a^2+5b^2=2$.
Irreducibility of $3$ was proved above. Each of the last two factors has
norm $6$; a nonunit split would require norms $2$ and $3$, both
impossible. The only units are $\pm1$, so no factor $2$ or $3$ is
associate to a factor with nonzero $\sqrt{-5}$ coordinate.
This is a domain but not a UFD. Merely writing two different products
would not suffice without these irreducibility and nonassociation checks.

\KEntry{K07.10.ufd_not_pid}{A UFD which is not a PID}
The polynomial theorem proved below in
Theorem~\ref{thm:polynomial-ufd} makes $\Z[X]$ a UFD.
Its ideal $(2,X)$ is proper, since its elements have even constant
coefficient. If generated by $h$, the divisibility $h\mid2$ forces
$h$ constant. Its integer value divides $2$, so it is $\pm1$ or $\pm2$.
The first alternatives generate the whole ring; the latter cannot divide
$X$, whose coefficient of $X$ is $1$. This contradiction proves that
$(2,X)$ is not principal.

% Topic grouping: A PID which is not Euclidean
\KEntry{K07.10.pid_not_euclidean}{The ring of discriminant minus nineteen}
Let $\omega=(1+\sqrt{-19})/2$ and $R=\Z[\omega]$.
Then $\omega^2-\omega+5=0$ and
\[
 N(a+b\omega)=(a+b\omega)(a+b\bar\omega)=a^2+ab+5b^2
 =(a+b/2)^2+19b^2/4 .
\]
This is a multiplicative, positive integer norm on nonzero elements.
If $b\ne0$, its value is at least $5$; hence its only value-$1$
elements, and its only units, are $\pm1$.

\begin{lemma}\label{lem:dedekind-hasse-sufficient}
Suppose a domain has a positive integer-valued norm $N$ such that,
whenever $a,b\ne0$ and $b\nmid a$, some $s,t\in R$ satisfy
$0<N(sa-tb)<N(b)$. Then it is a PID.
\end{lemma}
\begin{proof}
In a nonzero ideal choose $b\ne0$ of minimum positive norm.
For $a$ in the ideal not divisible by $b$, the asserted $sa-tb$ is
a nonzero element of smaller norm in that ideal, a contradiction.
Thus the ideal equals $(b)$.
\end{proof}
We verify this criterion, including the small denominators. For
$z=a_0/b_0\notin R$ in the fraction field write
\[
 z=\frac{a+b\sqrt{-19}}{c},\qquad
 a,b,c\in\Z,\ c>0,\ \gcd(a,b,c)=1 .
\]
Such a representation follows by rationalizing and clearing rational
denominators, then dividing their common gcd. We have $c>1$, since
$\Z[\sqrt{-19}]\subset R$. It suffices to find $s,t\in R$ with
$0<N(zs-t)<1$, because multiplication by $b_0$ gives the criterion.

If $c\ge5$, choose $x,y,h\in\Z$ with $ax+by+ch=1$.
Write $ay-19bx=cq+r$, with $|r|\le c/2$, taking a nearest multiple.
Set $s=y+x\sqrt{-19}$, $t=q-h\sqrt{-19}$. Then
$zs-t=(r+\sqrt{-19})/c$, and its norm is $(r^2+19)/c^2$.
For $c\ge6$ this is at most $1/4+19/36<1$; for $c=5$,
integrality gives $|r|\le2$, so it is at most $23/25<1$.
It is nonzero.
If $c=2$, $a,b$ have opposite parity: both even contradict the gcd,
and both odd would give $z\in R$. With $s=1$ and
$t=((a-1)+b\sqrt{-19})/2\in R$, the difference is $1/2$.
If $c=3$, the number $A=a^2+19b^2$ is nonzero modulo $3$,
since $a^2+b^2\equiv0\pmod3$ forces both divisible by $3$.
Write $A=3q+r$, $r=1,2$, and take
$s=a-b\sqrt{-19}$, $t=q$. The difference is $r/3$.
If $c=4$ and $a,b$ have opposite parity, $A$ is odd; write
$A=4q+r$, $r=1,3$, and use the same $s,t$, giving $r/4$.
If both are odd, $A\equiv4\pmod8$; write $A=8q+4$ and use
$s=(a-b\sqrt{-19})/2\in R$, $t=q$, giving $1/2$.
Both even are impossible. Each listed nonzero real difference has norm
strictly between zero and one. Thus $R$ is a PID.

To exclude every Euclidean structure, call a nonzero nonunit $u$ a
universal side divisor if for every $x\in R$, $u\mid x-\varepsilon$
for some $\varepsilon$ equal to zero or a unit.
In any Euclidean nonfield, a minimal nonzero nonunit $u$ has this
property: dividing $x$ by $u$ leaves zero or a smaller element, which
by minimality must be a unit.
In our ring, apply the property to $x=2$. Since the units are $\pm1$,
$u$ divides $1,2$, or $3$, and cannot divide $1$. Multiplicativity of
the norm forces $N(u)$ to divide $4$ or $9$.
There are no norm $2$ or $3$ elements, so $N(u)=4$ or $9$.
Norm $4$ forces $u=\pm2$. For norm $9$, $|b|\ge2$ is impossible;
for $b=\pm1$ the equation $a^2\pm a+5=9$ has no integral solution
(it would require the nonsquare discriminant $17$).
Thus norm $9$ forces $u=\pm3$.
But none of $\omega,\omega-1,\omega+1$ is divisible by $2$ or by $3$:
its $\omega$ coefficient is $1$. Taking $x=\omega$ therefore
contradicts the universal-side-divisor property.
The ring is not Euclidean for any well-founded Euclidean function,
not merely for its displayed norm.
